\chapter{Conclusioni}

\section{Obiettivi Raggiunti e Risultati Ottenuti}
Nel percorso di questa Tesi, si è analizzato, implementato e documentato un completo \textit{Migration Tool}
(con pacchetto aziendale \texttt{txamg}) nell'ambiente \textit{custom (tx)} Infor LN. I requisiti complessi, nati per la
maggior parte dalla natura multitenant dei nuovi ambienti CloudSuite, sono stati affrontati offrendo un approccio parametrico
e automatizzato.

Mediante lo sviluppo delle DLL di generazione del codice dal dizionario (ad esempio \texttt{txamgdll0012} e \texttt{txamgdll0013}) e
dell'architettura per la lettura nativa Infor Exchange, si è raggiunto l'obiettivo aziendale in termini di affidabilità,
performance logica e generazione del codice in tempo reale, eliminando la componente di trascrizione metodica manuale
e suscettibile d'errori.

\section{Sfide Incontrate}
Le principali complessità progettuali erano dovute al mapping rigido imposto dal requisiti interni e dall'interfacciarsi
ad entità Infor chiuse (\textit{closed source} o proprietari), in cui il Baan C (3GL e 4GL) risulta l'unica tecnologia
utilizzabile valida.

In un ecosistema di questa genere, l'accesso al compilatore interno garantendo una perfetta validazione del codice
per ogni variante di tenant è risultato cruciale dal punto di vista dell'analisi e debug dei prototipi iniziali della libreria generatrice di log.
Analogamente, standardizzare codifiche e formati non standard proveniente dagli strati Legacy ha richiesto robusti
componenti di parsing e una logica complessa di \textit{Runs} di batch (\texttt{txamg200}).

Durante lo fase di studio e sviluppo, è stato necessario affrontare una fase di training sia base che avanzata dei tool
Infor, quali struttura, nomenclatura, linguaggi, logica di funzionamento, architettura come \textit{middleware ION}, BOD, API e \textit{Web Services}
con formati XML legacy come SOAP (non REST), in tempi molto brevi rispetto a quanto normalmente previsto (diversi mesi), al fine
di completare il progetto in tempo e completare questa tesi.

Si fa anche notare la natura molto restrittiva dell'ERP Infor, che non permette l'accesso nemmeno ad amministratori di sistema
a parti di API o funzionalità che semplificherebbero lo sviluppo, oltre alla mancata documentazione delle stesse in molti casi
a causa di transizioni a più alto livello aziendale, lasciando allo sviluppatore finale (noi) "scoprire" ed esplorare
tramite trial-and-error delle soluzioni, spesso non standard, ove possibile.
Questo ha dato origine molti tempi morti, complicando il contesto generale


\section{Sviluppi Futuri}
Nonostante lo standard d'eccellenza raggiunto, si ritiene possibile evolvere ulteriormente il tool in diverse direzioni.

In primo luogo, lo schema architetturale per le estensioni \textit{Custom} potrebbe venire affiancato da procedure di \textit{Machine Learning} (addestrate sui mapping legacy)
in grado di suggerire i mapping e i moduli di conversione per grandi volumi di tabelle non ancora definite
(senza l'input umano in interfacce e Sessioni \texttt{txamg2103}).
Inoltre, la visualizzazione delle specifiche tramite esportazione
a modelli UML o grafi digitali, se automatizzata nativamente dal framework, fornirebbe reportistica in \textit{real-time} durante
la fase di studio del cliente prima dell'operazione di caricamento, riducendo di fatto tempi e costi
e supportando i consulenti AlteaIN, funzionali e Infor verso transizioni Cloud più semplici.
